\section{Introduction}

\label{sec:introduction}

% PAPER_STUDIO_PARAGRAPH:I-P1

Large language models are increasingly used for multiple-choice question answering, yet their responses may depend on the order in which answer options are presented. A model that reliably selects the correct answer under one ordering could fail under another, even when the semantic content is unchanged. This option-order sensitivity is consequential because it undermines the trustworthiness of models deployed in high-stakes settings, where users expect a consistent answer regardless of surface-level formatting. Moreover, such sensitivity complicates evaluation: benchmark scores that are computed under a single canonical ordering may overstate a model's true competence. The concrete problem addressed in this study is therefore to determine the extent to which random permutations of answer options alter model predictions, and to measure the consistency of model behavior across these orderings, as illustrated in Figure~\ref{fig:motivation}.

\begin{figure}[t]
  \centering
  \fbox{\parbox[c][6em][c]{0.92\linewidth}{\centering F1 placeholder: complete the final artwork after project export}}
  \caption{Large language models give different multiple-choice answers when option order changes; the figure poses the study's evaluation question of measuring prediction consistency across random permutations of answer choices.}
  \label{fig:motivation}
\end{figure}

% PAPER_STUDIO_PARAGRAPH:I-P2

Despite the practical importance of option-order sensitivity, existing work does not settle whether it reflects a systematic flaw or incidental variation \cite{}. Prior studies of prompt sensitivity have largely focused on instruction wording or example selection, leaving the specific effect of answer-option permutation underexplored \cite{}. Analyses that do examine ordering often treat it as a nuisance variable to be controlled rather than as an object of study in its own right. Consequently, there is no established account of how prevalent order-induced answer changes are across models, nor of whether certain models or question types are disproportionately affected. Without such evidence, it remains unclear whether current evaluation practices, which typically fix a single ordering, produce stable estimates of model competence or conceal systematic inconsistencies that a different ordering would reveal.

% PAPER_STUDIO_PARAGRAPH:I-P3

To address this gap, we propose a systematic study of permutation invariance in multiple-choice question answering. Our approach evaluates whether models produce consistent answers when the order of answer options is randomly permuted, treating the full set of permutations as the object of analysis rather than a nuisance factor. We define permutation invariance operationally as the stability of a model's selected answer across repeated orderings of the same question. The contributions of this study are threefold. First, we introduce a measurement protocol for quantifying order-induced answer variation. Second, we apply this protocol across a range of models to characterize the prevalence and magnitude of sensitivity. Third, we examine whether specific model families or question types exhibit differential susceptibility, providing a basis for understanding the underlying causes and for designing more robust evaluation practices.

% PAPER_STUDIO_PARAGRAPH:I-P4

The central evaluation question of this study is whether permutation invariance holds for large language models on multiple-choice question answering, and if not, how severe the violations are. Specifically, we ask: does randomly reordering answer options change a model's selected answer, and if so, how often and under what conditions? The claim boundary is deliberately narrow. We do not assert that permutation sensitivity, if observed, stems from any particular mechanism, such as positional bias, token-level confounds, or training artifacts. Instead, the study is designed to establish the empirical facts: the prevalence of answer changes across permutations, the consistency of model behavior under repeated orderings, and the extent to which sensitivity varies across models and question types. Within this boundary, the study aims to provide a reproducible measurement of permutation invariance that can inform evaluation design, without overclaiming causal explanations.
